\chapter{Discussion}
\label{ch:discussion}

% ===========================================================================
% WARNING, READ BEFORE WRITING. The original skeleton of this chapter asserted
% "LLM-agent coordination degrades more gracefully than pre-scripted
% workflows". THAT IS THE OPPOSITE OF WHAT HAPPENED. It was the proposal's
% expected finding, written in before the data existed. It has been removed.
% If any draft still carries it, delete it.
% ===========================================================================

\section{Interpretation}
\label{sec:disc:interpretation}
% Answer RQ1 and RQ2 directly, in operational terms, then stop.
%
% RQ1: No. The baseline saved every resident in every novel run; the best
% agentic configuration saved 90 %. Gate 1 fails 26 of 27 cells.
% Sharpen it with the regret result: the agentic system did not merely lose to
% a better system, it lost to INACTION. 68 of 150 runs came out worse than
% doing nothing.
%
% BUT SAY WHAT THE BASELINE'S WIN IS AND IS NOT. Its S = 1.000 is invariance,
% not recovery. It gained nothing by responding: regret exactly zero in all 50
% runs. Never write that it "adapted" or "replanned better"
% (Chapter~\ref{ch:critique}).
%
% RQ2: The deficit is a COMMITMENT failure, not a comprehension failure.
% Walk the eliminated alternatives, because the elimination is what makes the
% claim credible: wrong population (no, 150/150 correct), message layer (no,
% 93.8 % closure), too few missions or a withheld reserve (no, ~2.3 dispatches
% either way, bus-2 used in 100 % of runs), scorer dedup (no), driver
% over-claiming (no), running out of clock (mostly no).
% Then the positive evidence: the loss is QUANTISED at whole bus-loads, and
% plan revisions band monotonically against delivery.
%
% The rubric is where "they can reason, they cannot finish" is visible:
% parity on the five criteria about DECIDING, short on five about CLOSING THE
% LOOP. Do not round that to "decision-quality parity" - it is a split.

\section{The Capability Ordering}
\label{sec:disc:capability}
% Mean RD narrows with model capability: deepseek -0.27, glm-5.2 -0.16,
% qwen3.5 -0.10, consistent across all five novel disruptions.
% Under the regret reading this is "how much damage the response did", not
% "how far behind the baseline it fell". Say it that way.
% DO NOT EXTRAPOLATE past qwen3.5. The design does not support it, and there is
% no capability threshold observable within the final model set.
% qwen3.5 also strands far fewer residents than the others (2.0-5.0 per 100 vs
% 10-22.5), so its deficit has a materially different shape from deepseek's.

\section{Trade-offs}
\label{sec:disc:tradeoffs}
% The baseline wins response time, cost and determinism, and it also won the
% primary outcome. Report that without softening.
% The agentic cost is an order of magnitude: 9-10x messages per resident
% delivered, ~800k tokens, 1.8-2.7x wall-clock.
% The honest framing: this is a cost with no demonstrated benefit ON THESE
% SCENARIOS. Whether it ever buys anything is untested here, because no
% scenario in this thesis required it to.

\section{Transferability}
\label{sec:disc:transfer}
% Which domains share the two conditions this case was built around: a
% disruption space too large to pre-enumerate, and heterogeneous actors holding
% local information.
% CAUTION. This thesis did not demonstrate an agentic advantage under those
% conditions. What transfers is the METHOD and the DIAGNOSIS, not a positive
% result. Generalise the design considerations
% (Chapter~\ref{ch:design}), not the architecture's performance.

\section{Threats to Validity}
\label{sec:disc:validity}
% Construct, internal, external, conclusion validity.
% Construct validity of the primary DV is argued in
% Section~\ref{sec:eval:construct}; parity and the deviations are
% Chapter~\ref{ch:validity}; treatment validity is
% Chapter~\ref{ch:critique}. Summarise, do not repeat.

\section{Limitations}
\label{sec:disc:limits}
% State these plainly. Being caught not having noticed is worse than declaring.
% The full list with the position on each is in Chapter~\ref{ch:critique};
% keep this section short and point there.
%   - single case: one facility, one city, one hazard
%   - one SIM_TIME_SCALE (8.4), so H4 was not tested
%   - k = 10 (NOT the ~30 the proposal planned; Amendment 1 set it on budget
%     and precision grounds after power.py showed no k buys H1)
%   - the scenario-scaling sweep was never run
%   - no equity dimension: the resident population is homogeneous by
%     declaration, so a max-min subgroup difference is identically zero
%   - residents modelled as a passive count, not as agents
%   - the feasibility ceiling complicates cross-study comparison
%   - LLM nondeterminism; :cloud tags are not version-frozen
%   - the novel arm's third model changed four times, so that arm is
%     exploratory under PREREGISTRATION sec.0

\section{Lessons Learned}
\label{sec:disc:lessons}
% Short. What would be built differently, as a bridge into
% Chapter~\ref{ch:design}, which carries the generalisable version.
% Resist duplicating that chapter here.
